\documentclass[10pt,a4paper]{amsart}
\usepackage{amsmath,amssymb,amsthm,graphicx}
% Retained minimal preamble for the Stacks Project chapter more-algebra.tex (More on Algebra).
% The chapter file pinned-original-source.tex carries only \input{preamble} on its first line, so no
% preamble text lives inside the chapter itself.  The Stacks Project uses one single shared preamble for
% every chapter; the snapshot stages it identically as sources/stacks-chow/preamble.tex and
% sources/stacks-etale-cohomology/preamble.tex.  The block below is copied mechanically, line for line,
% from sources/stacks-chow/preamble.tex lines 14-18 (the Xy-pic setup) and 154-217 (the native theorem
% declarations together with the mathematical macros).  Nothing was added, removed or redefined; the
% wrapper class and the amsmath/amssymb/amsthm/graphicx base packages are supplied by the assembler.
% Chapter: more-algebra.tex -- More on Algebra.
% For commutative diagrams we use Xy-pic
\usepackage[all]{xy}

% We use 2cell for 2-commutative diagrams.
\xyoption{2cell}

% Theorem environments.
%
\theoremstyle{plain}
\newtheorem{theorem}[subsection]{Theorem}
\newtheorem{proposition}[subsection]{Proposition}
\newtheorem{lemma}[subsection]{Lemma}

\theoremstyle{definition}
\newtheorem{definition}[subsection]{Definition}
\newtheorem{example}[subsection]{Example}
\newtheorem{exercise}[subsection]{Exercise}
\newtheorem{situation}[subsection]{Situation}

\theoremstyle{remark}
\newtheorem{remark}[subsection]{Remark}
\newtheorem{remarks}[subsection]{Remarks}

\numberwithin{equation}{subsection}

% Macros
%
\def\lim{\mathop{\mathrm{lim}}\nolimits}
\def\colim{\mathop{\mathrm{colim}}\nolimits}
\def\Spec{\mathop{\mathrm{Spec}}}
\def\Hom{\mathop{\mathrm{Hom}}\nolimits}
\def\Ext{\mathop{\mathrm{Ext}}\nolimits}
\def\SheafHom{\mathop{\mathcal{H}\!\mathit{om}}\nolimits}
\def\SheafExt{\mathop{\mathcal{E}\!\mathit{xt}}\nolimits}
\def\Sch{\mathit{Sch}}
\def\Mor{\mathop{\mathrm{Mor}}\nolimits}
\def\Ob{\mathop{\mathrm{Ob}}\nolimits}
\def\Sh{\mathop{\mathit{Sh}}\nolimits}
\def\NL{\mathop{N\!L}\nolimits}
\def\CH{\mathop{\mathrm{CH}}\nolimits}
\def\proetale{{pro\text{-}\acute{e}tale}}
\def\etale{{\acute{e}tale}}
\def\QCoh{\mathit{QCoh}}
\def\Ker{\mathop{\mathrm{Ker}}}
\def\Im{\mathop{\mathrm{Im}}}
\def\Coker{\mathop{\mathrm{Coker}}}
\def\Coim{\mathop{\mathrm{Coim}}}

% Boxtimes
%
\DeclareMathSymbol{\boxtimes}{\mathbin}{AMSa}{"02}

%
% Macros for moduli stacks/spaces
%
\def\QCohstack{\mathcal{QC}\!\mathit{oh}}
\def\Cohstack{\mathcal{C}\!\mathit{oh}}
\def\Spacesstack{\mathcal{S}\!\mathit{paces}}
\def\Quotfunctor{\mathrm{Quot}}
\def\Hilbfunctor{\mathrm{Hilb}}
\def\Curvesstack{\mathcal{C}\!\mathit{urves}}
\def\Polarizedstack{\mathcal{P}\!\mathit{olarized}}
\def\Complexesstack{\mathcal{C}\!\mathit{omplexes}}
% \Pic is the operator that assigns to X its picard group, usage \Pic(X)
% \Picardstack_{X/B} denotes the Picard stack of X over B
% \Picardfunctor_{X/B} denotes the Picard functor of X over B
\def\Pic{\mathop{\mathrm{Pic}}\nolimits}
\def\Picardstack{\mathcal{P}\!\mathit{ic}}
\def\Picardfunctor{\mathrm{Pic}}
\def\Deformationcategory{\mathcal{D}\!\mathit{ef}}

% --- appended 2026-10-12 (Chief): the project own definitions of the reference/slogan/history
% display environments, copied mechanically from the shared Stacks preamble lines 8-12
% (\usepackage{verbatim} plus the three \newenvironment lines, which the chapter preamble keeps
% verbatim from the project preamble). Nothing added, removed or redefined: these environments
% carry no mathematical content and the project itself typesets them as comments. Without these
% three lines a counted statement whose own text uses them fails with Environment slogan undefined.
\usepackage{verbatim}
\newenvironment{reference}{\comment}{\endcomment}
\newenvironment{slogan}{\comment}{\endcomment}
\newenvironment{history}{\comment}{\endcomment}
\usepackage{cleveref}
\crefname{originallemma}{original lemma}{original lemmas}
\Crefname{originallemma}{Original lemma}{Original lemmas}
\makeatletter
\newlabel{more-algebra-lemma-independence-of-generators}{{066A}{source}{}{sourcepointer-1}{}}
\newlabel{more-algebra-lemma-independence-of-generators@cref}{{[originallemma][0][]066A}{[1][0][]source}{}{}{}}
\makeatother
\crefname{originallemma}{original lemma}{original lemmas}
\Crefname{originallemma}{Original lemma}{Original lemmas}
\makeatletter
\newlabel{lemma-functorial}{{0624}{source}{}{sourcepointer-2}{}}
\newlabel{lemma-functorial@cref}{{[originallemma][0][]0624}{[1][0][]source}{}{}{}}
\makeatother
\crefname{originallemma}{original lemma}{original lemmas}
\Crefname{originallemma}{Original lemma}{Original lemmas}
\makeatletter
\newlabel{lemma-homotopy-koszul}{{0663}{source}{}{sourcepointer-3}{}}
\newlabel{lemma-homotopy-koszul@cref}{{[originallemma][0][]0663}{[1][0][]source}{}{}{}}
\makeatother
\crefname{originallemma}{original lemma}{original lemmas}
\Crefname{originallemma}{Original lemma}{Original lemmas}
\makeatletter
\newlabel{lemma-H1-regular-quasi-regular}{{062I}{source}{}{sourcepointer-4}{}}
\newlabel{lemma-H1-regular-quasi-regular@cref}{{[originallemma][0][]062I}{[1][0][]source}{}{}{}}
\makeatother
\title{Regularity of a generating sequence does not depend on the chosen generators of the ideal}
\author{The Stacks Project Authors}
\date{stacks-project:more-algebra.tex:local-snapshot-20261012; GNU FDL 1.2 or later}
\begin{document}
\maketitle
\setcounter{subsection}{1}
\noindent\textit{Editorial note.} All delivered mathematical text is the author's own, copied from the pinned Stacks Project chapter file more-algebra.tex; no statement or proof step was rewritten, reordered or removed.\par\medskip
\noindent\textit{Editorial note.} The chapter file pinned-original-source.tex carries only \textbackslash{}input\{preamble\} on its first line, so no preamble text lives inside the chapter. The retained file sources/stacks-more-algebra/retained-preamble.tex is a mechanical, line-for-line copy of the project's single shared preamble as staged in the snapshot (Xy-pic setup plus theorem declarations and mathematical macros). Nothing was added, removed or redefined, and the excerpt serves only as preamble, never as a span coordinate.\par\medskip
\noindent\textit{Editorial note.} Both counted spans start at the \textbackslash{}begin line of their environment, so the statement and the author's proof are delivered whole; surrounding prose of the section is deliberately outside the counted spans and is carried as standard prerequisites instead.\par\medskip
\noindent\textit{Editorial note.} The citation fields (source\_id, source\_version, authors, source\_url, license.name, source\_title) are identical to those of every other card this lane has produced from this chapter, so the batch selector treats them as one source.\par\medskip
\noindent\textit{Allowed standard prerequisites (editorial).}\begin{itemize}
\item Quasi-regularity, H\_1-regularity and Koszul-regularity of a sequence, and the equivalences between them for sequences generating the ideal. the author cites Lemma lemma-H1-regular-quasi-regular at proof lines 7740-7742. Supplies the reduction of the last two parts to the quasi-regular case.
\item The change-of-basis matrix between two generating sequences induces a map of Koszul complexes which is an isomorphism after inverting its determinant. the author cites Lemma lemma-functorial at proof lines 7730-7733. Produces the comparison map between the two Koszul complexes.
\item The cohomology modules of the Koszul complex of a sequence generating the ideal are annihilated by the ideal. the author cites Lemma lemma-homotopy-koszul at proof lines 7734-7736. Upgrades the local isomorphism to a quasi-isomorphism over the whole ring.
\end{itemize}\medskip
\ifdefined\hypertarget\hypertarget{sourcepointer-1}{}\fi
\noindent\textit{Source reference.} Original lemma 066A: focused-seeds/AGP/sources/stacks-snapshot-20261011/tags.txt, source lines 7929--7929. The lemma's own canonical tag, quoted so that readers can resolve the label inside the delivered text; it points at the counted statement itself and carries no proof content.\par\medskip
\ifdefined\hypertarget\hypertarget{sourcepointer-2}{}\fi
\noindent\textit{Source reference.} Original lemma 0624: focused-seeds/AGP/sources/stacks-snapshot-20261011/tags.txt, source lines 7783--7783. Navigation pointer to a statement the counted proof refers to but does not contain; nothing load-bearing is delegated to it.\par\medskip
\ifdefined\hypertarget\hypertarget{sourcepointer-3}{}\fi
\noindent\textit{Source reference.} Original lemma 0663: focused-seeds/AGP/sources/stacks-snapshot-20261011/tags.txt, source lines 7922--7922. Navigation pointer to a statement the counted proof refers to but does not contain; nothing load-bearing is delegated to it.\par\medskip
\ifdefined\hypertarget\hypertarget{sourcepointer-4}{}\fi
\noindent\textit{Source reference.} Original lemma 062I: focused-seeds/AGP/sources/stacks-snapshot-20261011/tags.txt, source lines 7797--7797. Navigation pointer to a statement the counted proof refers to but does not contain; nothing load-bearing is delegated to it.\par\medskip
\begin{lemma}
\label{lemma-independence-of-generators}
Let $R$ be a ring. Let $I$ be an ideal generated by $f_1, \ldots, f_r \in R$.
\begin{enumerate}
\item If $I$ can be generated by a quasi-regular sequence of length $r$,
then $f_1, \ldots, f_r$ is a quasi-regular sequence.
\item If $I$ can be generated by an $H_1$-regular sequence of length $r$,
then $f_1, \ldots, f_r$ is an $H_1$-regular sequence.
\item If $I$ can be generated by a Koszul-regular sequence of length $r$,
then $f_1, \ldots, f_r$ is a Koszul-regular sequence.
\end{enumerate}
\end{lemma}

\begin{proof}
If $I$ can be generated by a quasi-regular sequence of length $r$,
then $I/I^2$ is free of rank $r$ over $R/I$. Since $f_1, \ldots, f_r$
generate by assumption we see that the images $\overline{f}_i$ form a basis of
$I/I^2$ over $R/I$. It follows that $f_1, \ldots, f_r$ is a quasi-regular
sequence as all this means, besides the freeness of $I/I^2$, is that the maps
$\text{Sym}^n_{R/I}(I/I^2) \to I^n/I^{n + 1}$ are isomorphisms.

\medskip\noindent
We continue to assume that $I$ can be generated by a
quasi-regular sequence, say
$g_1, \ldots, g_r$. Write $g_j = \sum a_{ij}f_i$. As $f_1, \ldots, f_r$
is quasi-regular according to the previous paragraph, we see that
$\det(a_{ij})$ is invertible mod $I$. The matrix
$a_{ij}$ gives a map $R^{\oplus r} \to R^{\oplus r}$ which induces
a map of Koszul complexes
$\alpha : K_\bullet(R, f_1, \ldots, f_r) \to K_\bullet(R, g_1, \ldots, g_r)$,
see
Lemma \ref{lemma-functorial}.
This map becomes an isomorphism on inverting $\det(a_{ij})$.
Since the cohomology modules of both $K_\bullet(R, f_1, \ldots, f_r)$ and
$K_\bullet(R, g_1, \ldots, g_r)$ are annihilated by $I$, see
Lemma \ref{lemma-homotopy-koszul},
we see that $\alpha$ is a quasi-isomorphism.

\medskip\noindent
Now assume that $g_1, \ldots, g_r$ is a $H_1$-regular sequence generating $I$.
Then $g_1, \ldots, g_r$ is a quasi-regular sequence by
Lemma \ref{lemma-H1-regular-quasi-regular}. By the previous paragraph
we conclude that $f_1, \ldots, f_r$ is a $H_1$-regular sequence.
Similarly for Koszul-regular sequences.
\end{proof}

\end{document}
